%% kapa: exam=1 round=24 year=2013 period=1 form=A subject=civil_law subject_ko=민법 seq=3 number=3
%% source: Q-Net 감정평가사 기출문제 (cst00309) artlSeq=5073802 file=extracted/1008/5073802/2013년 제24회 감정평가사 1차시험 기출문제/1교시(A형).hwp
%% conversion: kordoc 4.24.51 -> tools/gen1.py (auto); review=proofread-sonnet; flags=none
\begin{question}{3}
미성년자의 행위능력에 관한 설명으로 옳지 않은 것은? (다툼이 있으면 판례에 의함)
\begin{choices}
  \item 미성년자에게 영업을 허락한 법정대리인은 미성년자의 동의가 없더라도 그 영업허락을 제한할 수 있다.
  \item 법정대리인이 영업허락을 취소했으나 그 사실을 모르는 제3자가 미성년자와 체결한 영업상의 계약은 유효하다.
  \item 의사능력 있는 미성년자가 타인으로부터 대리권을 수여받아 부모의 동의 없이 대리행위를 한 경우 무능력을 이유로 그 대리행위를 취소할 수 없다.
  \item 미성년자가 사술(詐術)로 자기를 성년자로 믿게 했더라도 미성년을 이유로 그 행위를 취소할 수 있다.
  \item 미성년자의 법률행위에 법정대리인의 동의가 있었는지에 대한 증명책임은 동의 있음을 주장하는 상대방에게 있다.
\end{choices}
\end{question}
\answer{4}
